\section{Requirements Overview}
\label{sec:requirements}

This section formalises the functional and non-functional requirements that the architecture described in Section~\ref{sec:architecture} is designed to satisfy. Requirements are derived from the project goals in Section~\ref{sec:app-introduction} and the operational needs of HPC administrators and researchers.

\subsection{Functional Requirements}

\begin{table}[H]
    \centering
    \caption{Functional Requirements}
    \label{tab:functional-requirements}
    \small
    \begin{tabular}{|l|p{10.5cm}|}
        \hline
        \textbf{ID} & \textbf{Requirement} \\
        \hline
        FR1 & The system shall interpret natural-language Slurm queries, including variations in phrasing and implicit context, and map them to appropriate typed tool calls. \\
        \hline
        FR2 & The system shall retrieve live cluster state (queues, nodes, partitions, accounting, reservations, scheduler diagnostics) via read-only Slurm tools and present results grounded in tool output. \\
        \hline
        FR3 & The system shall execute state-changing operations (submit, cancel, hold, release, requeue, node control, account management) only after explicit user confirmation. \\
        \hline
        FR4 & The system shall retrieve relevant Slurm documentation from a local corpus using hybrid retrieval (lexical + semantic), with web search as a fallback when the local corpus does not cover the topic. \\
        \hline
        FR5 & The system shall maintain multi-turn conversation context, enabling follow-up questions and contextual references within a session. \\
        \hline
        FR6 & The system shall provide real-time feedback during execution through structured status events, making tool selections and intermediate results visible to the user. \\
        \hline
        FR7 & The system shall support deterministic evaluation against mock cluster states, producing behavioral traces that can be compared to expected tool selections, routing decisions, and state transitions. \\
        \hline
    \end{tabular}
\end{table}

\subsection{Non-Functional Requirements}

\begin{table}[H]
    \centering
    \caption{Non-Functional Requirements}
    \label{tab:nonfunctional-requirements}
    \small
    \begin{tabular}{|l|p{10.5cm}|}
        \hline
        \textbf{ID} & \textbf{Requirement} \\
        \hline
        NFR1 & \textbf{Safety}: The system shall never execute cluster-mutating commands without explicit confirmation. The confirmation mechanism shall be enforced at the tool layer, independent of LLM behavior. \\
        \hline
        NFR2 & \textbf{Accuracy}: Responses shall be grounded in tool output. The system shall not fabricate cluster state, job IDs, node names, or resource figures. \\
        \hline
        NFR3 & \textbf{Latency}: Simple read-only queries shall complete within interactive response times. Long-running workflows shall emit incremental status events. \\
        \hline
        NFR4 & \textbf{Extensibility}: New Slurm tools shall be addable via the MCP server without modifying the agent core. Tool schemas are discovered at runtime. \\
        \hline
        NFR5 & \textbf{Context Efficiency}: The system shall avoid unnecessary growth of LLM context by lazy-loading documentation, capping retrieval output size, and separating large artifacts from conversational memory. \\
        \hline
        NFR6 & \textbf{Reproducibility}: Evaluation runs shall be deterministic given the same mock state, model, and test case. Provider and model shall be configurable per session. \\
        \hline
    \end{tabular}
\end{table}
